% Содержимое подраздела 2.2 - Обучение с подкреплением (Reinforcement Learning) для прямой оптимизации свойств


\section{Обучение с подкреплением (Reinforcement Learning) для прямой оптимизации свойств}

Подходы, рассмотренные в предыдущем разделе — вариационные автоэнкодеры и гетероэнкодеры — реализуют непрямой контроль над генерацией. Их общая стратегия заключается в построении качественного латентного пространства и последующей оптимизации внутри него.

Принципиально иной парадигмой целевой генерации является обучение с подкреплением (Reinforcement Learning, RL), которое реализует прямой контроль над свойствами. В отличие от подходов, основанных на оптимизации в латентном пространстве (таких как VAE и гетероэнкодеры), RL-методы не пытаются построить глобальную модель химического пространства, а напрямую обучают генеративную модель (агента) максимизировать целевую функцию (вознаграждение) путем итеративного взаимодействия со «средой» (вычислительной моделью, оценивающей свойства).

\rgnsubsection{Теоретические основы: фреймворк REINVENT}

Одним из наиболее известных и широко используемых фреймворков для RL-дизайна молекул является REINVENT \cite{loeffler2024reinvent}. Его работа основана на взаимодействии трех ключевых компонентов:
\begin{itemize}
    \item \textbf{Агент (Agent):} Генеративная модель, как правило, рекуррентная нейронная сеть (RNN) на основе LSTM или Трансформер, которая обучена генерировать молекулы в виде SMILES-строк.
    \item \textbf{Среда (Environment):} Программный модуль, который принимает на вход сгенерированную молекулу и вычисляет для нее скалярную оценку (вознаграждение) в диапазоне от 0 до 1. Среда может включать в себя предсказательные QSAR-модели, расчеты молекулярного докинга, оценку физико-химических свойств и т.д.
    \item \textbf{Вознаграждение (Reward):} Целевая функция, которую агент стремится максимизировать.
\end{itemize}
Центральным механизмом обучения является оптимизация политики агента с помощью стратегии "Difference between Augmented and Posterior" (DAP). Функция потерь минимизирует расхождение между правдоподобием молекулы $T$ у текущего агента ($\log P_{\text{agent}}(T)$) и так называемым <<дополненным>> правдоподобием:
$$ \log P_{\text{aug}}(T) = \log P_{\text{prior}}(T) + \sigma S(T) $$
где $P_{\text{prior}}(T)$ -- правдоподобие от исходной, предварительно обученной модели (prior), $S(T)$ -- итоговая оценка (вознаграждение) молекулы, а $\sigma$ -- гиперпараметр, контролирующий силу оптимизации. Такой подход позволяет агенту генерировать высокооцененные молекулы, не слишком отдаляясь от химически <<правдоподобного>> пространства, изученного prior-моделью, что предотвращает катастрофическое забывание.

\rgnsubsection{Пример адаптации: REINVENT с GPU-ускорением для поиска ингибиторов ВИЧ-1}

Стандартные реализации RL-пайплайнов часто сталкиваются с вычислительным <<бутылочным горлышком>>, когда в качестве среды используется ресурсоемкий молекулярный докинг на CPU. Для решения этой проблемы была предложена адаптация архитектуры REINVENT, в которой стандартный механизм оценки был заменен на высокопроизводительный конвейер с использованием докинга на GPU \cite{воробьев2024адаптация}.

Ключевой модификацией стала интеграция программы \\ `AutoDock-Vina-GPU-2.1` непосредственно в цикл обучения RL. На каждом шаге модель генерирует пакет SMILES-строк, которые в реальном времени конвертируются в 3D-структуры и отправляются на GPU-докинг в активный сайт белка-мишени gp120 ВИЧ-1. Для повышения точности предсказаний результаты докинга дополнительно переоценивались с помощью модели машинного обучения `RFScore-4`, которая показала значительно лучшую корреляцию с экспериментальными данными (коэффициент Пирсона 0.661 против 0.476 у Vina-GPU).

Функция вознаграждения также была усложнена. Базовое вознаграждение, вычисленное из энергии связывания $\Delta G$, умножалось на штрафные коэффициенты в случае нарушения молекулой ключевых физико-химических свойств (например, правила пяти Липинского). Такой подход направляет генерацию в сторону не просто активных, но и биодоступных соединений.
В результате за 60 часов работы на GPU было сгенерировано более 62 000 уникальных молекул, из которых более 34 000 удовлетворили как критерию высокого сродства ($\Delta G < -9.2$ ккал/моль), так и всем фармакокинетическим ограничениям, став перспективными кандидатами для дальнейших исследований \cite{воробьев2024адаптация}.

\rgnsubsection{Повышение эффективности RL: Curriculum Learning и Sample Efficiency}

Несмотря на свою гибкость, policy-based RL страдает от двух существенных проблем: неэффективности обучения на сложных задачах и низкой выборочной эффективности при работе с дорогими оракулами.

Для решения первой проблемы был предложен подход \textbf{Обучения по Учебному Плану (Curriculum Learning, CL)} \cite{guo2022improving}. Вместо того чтобы сразу решать сложную многопараметрическую задачу, агент сначала обучается на последовательности упрощенных задач. Например, при оптимизации с использованием докинга, на первом этапе агент может обучаться генерировать молекулы, просто похожие по 2D-структуре на известный лиганд (быстрая оценка), и только после достижения определенного порога в эту задачу добавляется дорогостоящий докинг. Эксперименты показали, что стандартный RL-агент не генерировал продуктивных молекул в течение первых 150-200 эпох, в то время как CL-подход начал делать это практически немедленно, сгенерировав в итоге в 3.7-4.2 раза больше целевых соединений.

Вторую проблему -- низкую \textbf{выборочную эффективность (sample inefficiency)} -- решает алгоритм \textbf{Augmented Memory} \cite{guo2024augmented}. Он радикально повышает эффективность использования каждого дорогостоящего вызова оракула за счет комбинации техник <<повторения опыта>> (experience replay) и <<аугментации данных>> (SMILES augmentation). Все лучшие найденные молекулы сохраняются в буфере памяти. На каждом шаге обучения агент многократно обновляется на основе всех молекул из этого буфера, причем для каждой молекулы генерируется множество ее эквивалентных SMILES-представлений. Это создает чрезвычайно сильный обучающий сигнал. Данный метод установил новый state-of-the-art на стандартизированном бенчмарке PMO, достигнув AUC Top-10 = 15.002 (против 14.016 у базового REINVENT) и показав способность генерировать в 2-3 раза больше качественных молекул при том же бюджете вызовов оракула.
